\section{RQ7: Scope across Methods, Scale, and Graph Family}
\label{sec:families}

\paragraph{External adaptive methods.}
The frozen DARTH adapters produce target query-failure rates of 22.55\% on SIFT-100K and 23.12\% on Arxiv-Nomic-100K at $\rec@10<0.95$. Ada-ef yields 9.52 [9.44, 9.60]\% on ten Arxiv targets, with a target-build bootstrap conditional on its query set. All ten Ada-ef candidate certificates fail; audited execution uses the qualified endpoint and has zero search gain relative to that endpoint. Ada-ef is absent from Euclidean SIFT because the inspected estimator implements cosine/inner-product semantics, not a semantics-preserving L2 path. Its exclusion is an implementation-scope boundary, not a failed SIFT result.

These measurements show why query-risk qualification matters beyond TCP. They do not, without matched source checks, identify how much failure was caused by rebuilding, and the event is not interchangeable with mean recall or expected false-negative fraction. They therefore remain qualification studies rather than a ranking of published methods. The TCP comparator and information budget are stated separately.

\paragraph{Larger scale.}
On Deep1M, eight hnswlib builds cross four seeds with two insertion orders, giving 56 directed pairs on 1,000 queries. Source-derived first-passing actions have 28.86\% transport risk against 8.82\% target-reference risk. The incremental risk is 20.04 [17.29, 22.77] percentage points, using a target-build bootstrap conditional on the frozen source histories and query set. Keeping the reference risk visible matters: this is additional loss above a difficult, partly unresolved target reference, not a transfer from a zero-risk baseline.

A preregistered replication reuses the eight indexes with disjoint 500/500/500-query roles. All 56 non-clipped candidates pass target certification; evaluation risk is 1.63 [0.80, 2.61]\% and mean-NDC saving versus efSearch 1200 is 34.28 [33.53, 35.05]\%. Every target and leave-one-target-out analysis remains positive; this is scale, not cross-family, confirmation.

\paragraph{Vamana-style evidence.}
Six source and six separate target builds produce 36 directions and 750 queries per dataset, using native search-list size and beam width one. The frozen event table combines its stage-specific under-budget and censoring outcomes. A target-build reanalysis gives 16.86 [15.88, 17.93]\% on SIFT and 11.37 [10.80, 11.99]\% on Arxiv. Its recorded responses do not reconstruct exactly the first-passing execution estimand used for hnswlib/Faiss. We retain it as directional cross-family evidence, not another row of Table~\ref{tab:graph-only}.

Vamana cost-tax intervals include zero: 1.07 [$-1.03$, 2.86]\% and 1.17 [$-0.83$, 3.03]\%, respectively, for the mean of per-target ratios in that reanalysis. Thus it supplies no established positive cost penalty or cross-engine superiority. Table~\ref{tab:scope} summarizes which question each setting can answer.

\begin{table*}[t]
\caption{Scope of the extension evidence. Each row retains its event, reference, and inference unit. All numeric intervals in the text are 95\%; no common cross-engine ranking or pooled effect is formed.}
\label{tab:scope}
\small\begin{tabularx}{\textwidth}{@{}lXXl@{}}\toprule
Setting & Event and reference & Supported comparison & Unit / query count\\\midrule
DARTH & Absolute target $\rec@10<0.95$ risk; no matched reference subtraction & Target qualification under the frozen adapter & Target summaries; 1,000/build\\
Ada-ef & Absolute target $\rec@10<0.95$ risk; endpoint for audited saving & Candidate rejection versus endpoint execution on Arxiv & 10 target builds; 1,000/build\\
Deep1M & Transport versus target reference; certified slack versus efSearch 1200 & Transfer loss and certified recovery at 1M scale & 8 targets; 1,000 prior + 1,500 fresh\\
Vamana-style & Frozen stage under-budget/censoring event; stage-specific cost reference & Directional cross-family transfer evidence, not a harmonized HNSW increment & 6 target builds; 750 shared queries\\\bottomrule
\end{tabularx}
\end{table*}

Together with RQ1, these results support a layered conclusion: response non-portability is seen across HNSW implementations and at larger scale; qualification of adaptive actions is operationally consequential; and a distinct graph family supplies supporting evidence with narrower semantic comparability. They do not establish uniform failure across every graph family or every adaptive policy.
